\documentclass[aspectratio=169, 12pt]{beamer}
%\usetheme{Madrid}
\usetheme{default}
\usecolortheme{default}

\title{Achieving Consensus on a Blockchain}

\author{Thomas Sargent \\
World Laureates Association \\
Dubai, UAE}
\date{February 2026}
%\institute{World Laureates Association}

\begin{document}
\titlepage
%% Title slide
%\begin{frame}
%    \titlepage
%    \centering
%    \small Game-theoretic analysis of blockchain consensus, \\
%    checkpoint strategies, and the efficiency-consensus trade-off.
%\end{frame}

\begin{frame}{Important Paper}


\textbf{\em Game-theoretic analysis of blockchain consensus,
    checkpoint strategies, and  efficiency-consensus trade-off,
         Zetlin-Jones et al. (2025)}
         
\end{frame}

\begin{frame}{Game-theoretic model of blockchain consensus with rational miners }
  %  \textbf{Key contributions}:
    \begin{enumerate}
        %\item Game-theoretic model of blockchain consensus with rational miners.
        \item Recognizes three attack types under longest-chain rule.
        \item \textbf{Checkpoint strategies} as history-dependent solution:
        \begin{itemize}
            \item Prevent double-spending.
            \item Enhance consensus robustness.
        \end{itemize}
        \item Parent-block checkpoints can be used to manage latency.
        \item There is a  \textbf{Grossman-Stiglitz-like trade-off}  between informational efficiency and consensus.
            \end{enumerate}
    
    \vspace{0.3cm}
    \textbf{Implications}: Presence of vulnerabilities is endogenous and  depends on equilibrium strategies. Modifications of protocols (some already in use) can significantly improve security.
\end{frame}



\begin{frame}{Object in Play}        

\begin{itemize}
        \item \textbf{Blockchain}: a decentralized ledger managed  without trusted third parties.
        \begin{itemize}
        \item a directed graph (arborescence); each path = possible ledger.
        \end{itemize}
\end{itemize}

\end{frame}



% Slide 1: Introduction and Motivation
\begin{frame}{The Difficulty}
    \begin{itemize}
        \item \textbf{Blockchains}: decentralized ledgers managed communally by selfish transaction verifiers, e.g.., Bitcoin``miners''.
        \item \textbf{Key challenge}: achieving \textbf{consensus} among self-interested users.
        \item Bitcoin’s \textbf{longest-chain rule} treats miners as rational agents.
        \item Prior work shows it fails under certain conditions:
        \begin{itemize}
            \item Biais et al. (2019), Budish (2025): double-spending possible.
            \item Vulnerable when miners are heterogeneous in several senses.
        \end{itemize}
        \item \textbf{ Zetlin-Jones et al. (2026)}: construct  a \textbf{dynamic game-theoretic model} to:
        \begin{itemize}
            \item Study consensus as an equilibrium outcome.
            \item Propose \textbf{checkpoint strategies} to improve robustness.
        \end{itemize}
    \end{itemize}
\end{frame}

% Slide 2: Model Overview
\begin{frame}{Model Environment}
    \begin{itemize}
        \item \textbf{Miners}: \(N\) infinite-lived, discount factor \(\delta\), mining power \(p_i\).
        \item \textbf{Time}: discrete periods, miners choose block location.
        \item \textbf{Blockchain}: directed graph (arborescence); each path = possible ledger.
        \item \textbf{Sources of Utility}:
        \begin{itemize}
            \item Coin balances (more valuable if more miners agree).
            \item Consumption goods (via spend transactions, settled with delay).
        \end{itemize}
        \item \textbf{Equilibrium concept}: Perfect Public Equilibrium (public strategies).
        \item \textbf{Ultimate Source of Consensus}:  balances can be  spent only if miners agree. %d upon, miners value consensus.
    \end{itemize}
\end{frame}


\begin{frame}
\frametitle{Dimensions of Miners' Heterogeneity}

\textbf{Authors' Innovation:} Specify miners' preferences \textbf{directly} over information recorded on different ledger versions

\vspace{0.5em}

\begin{enumerate}
    \item \textbf{Mining Power} $(p_i)$
    \begin{itemize}
        \item Probability miner $i$ successfully appends a block
        \item $\{p_i\}$'s are  exogenous;  $\sum_{i=1}^N p_i = 1$
        \item Can be highly concentrated in practice
    \end{itemize}

    \vspace{0.5em}

    \item \textbf{Coin Balances} $(Y_{i,b} + R_{i,b})$
    \begin{itemize}
        \item $R_{i,b} = \bar{R}$ if miner $i$ solved block $b$; zero otherwise
        \item $Y_{i,b} \in \mathbb{R}$: transaction data (positive = received, negative = spent)
        \item Balances vary across blocks and chains
    \end{itemize}
\end{enumerate}

\end{frame}

\begin{frame}
\frametitle{Heterogeneous Preferences Over Chains}

\textbf{Flow utility for miner $i$:}
\[
u^i(q_{i,t}; H_t) = \sum_{b \in \mathcal{B}(G_t)} \left[ (1-\delta) q_{i,b,t}(Y_{i,b} + R_{i,b}) + \frac{Y_{i,b}^-}{\delta} \cdot \lambda_t(b, H_t) \right]
\]

\vspace{0.5em}

\textbf{Key elements:}
\begin{itemize}
    \item $q_{i,b,t}$: weight on block $b$ = share of \emph{other} miners' power on chains containing $b$
    \item Balances more valuable when more miners ``agree'' (work on same chain)
    \item Spend transactions $(Y_{i,b}^-)$ yield consumption once settled
\end{itemize}

\vspace{0.5em}

\textbf{Implication:} Miners prefer \emph{different} chains
\[
b^*_{i,t} = \arg\max_{b \in \mathcal{B}^{LC}(G_t)} \sum_{b' \in \mathcal{C}(b, G_t)} \left( Y_{i,b'} + R_{i,b'} + \frac{Y_{i,b'}^-}{\delta} \cdot \Lambda_t(b', H_t) \right)
\]

\end{frame}
% Slide 3: Longest-Chain Rule – The Problem
\begin{frame}{Longest-Chain Rule has a Problem}
    The Rule: extend chain with most blocks.
    
    \vspace{0.3cm}
    \textbf{Fails under three profitable choices (deviations)}:
    \begin{enumerate}
        \item \textbf{Double-Spending Attacks}
        \begin{itemize}
            \item Fork to omit large spend transaction.
            \item Profitable if spend \(\gg\) block rewards (Budish 2025).
        \end{itemize}
        \item \textbf{Consensus Redirection Attacks}
        \begin{itemize}
            \item Fork to include favorable balances not on longest chain.
            \item Profitable with concentrated mining power.
        \end{itemize}
        \item \textbf{Saving-on-Consensus Attacks}
        \begin{itemize}
            \item Fork to delay spend inclusion, reduce flow disutility.
            \item Profitable for high-power miners.
        \end{itemize}
    \end{enumerate}
    
  \textbf{  \alert{With heterogeneous miners, longest-chain rule fails}}
\end{frame}

\begin{frame}
\frametitle{Heterogeneity Matters}

When miners are \textbf{sufficiently heterogeneous}, the longest chain rule fails:

\vspace{0.5em}

\begin{enumerate}
    \item \textbf{Double-Spending Attacks}
    \begin{itemize}
        \item Miner with large spend transaction $(Y_{i,b} < 0)$ on longest chain
       \begin{itemize} \item Incentive to fork and omit the spend $\Rightarrow$ ``undo'' the payment
        \item Profitable when: $(1 + p_i \delta)\bar{R} < p_i^2 \delta^2 |Y_{i,b}|$
        \end{itemize}
    \end{itemize}

    \vspace{0.3em}

    \item \textbf{Consensus Redirection Attacks}
    \begin{itemize}
        \item Miner with high $p_i$ and large balances on a shorter fork
        \begin{itemize}
        \item Incentive to extend fork rather than longest chain
        \end{itemize}
    \end{itemize}

    \vspace{0.3em}

    \item \textbf{Saving on Consensus Attacks}
    \begin{itemize}
        \item Miner with pending spend transaction in current block
        \begin{itemize}
        \item Incentive to delay consensus to reduce flow disutility
        \end{itemize}
    \end{itemize}
\end{enumerate}

\vspace{0.5em}

\textbf{Resolution:} Checkpoint strategies introduce history-dependence that aligns incentives

\end{frame}



% Slide 4: Checkpoint Strategies – A Solution
\begin{frame}{Possible Solution: Checkpoint Strategies }
    \textbf{Idea}: History-dependent rule – ignore forks that omit blocks behind a \textbf{checkpoint}.
    \begin{itemize}
        \item Checkpoint = reference block updated periodically.
        \item Miners consider only  chains containing the checkpoint.
        \item \textbf{Key property}: After a block is checkpointed, earlier transactions cannot be removed.
        \item Eliminates double-spending for checkpointed blocks (Proposition 3).
        \item Can sustain consensus even with:
        \begin{itemize}
            \item Large transactions.
            \item Concentrated mining power (Proposition 4).
        \end{itemize}
        \item \textbf{Practical use}: Ethereum, Polkadot  use checkpoint-like finality gadgets.
    \end{itemize}
\end{frame}

% Slide 5: Latency and Checkpoint Robustness
\begin{frame}{Latency and Checkpoint Robustness}
    \textbf{Network latency}: miners see different blocks temporarily.
    \begin{itemize}
        \item Can cause \textbf{checkpoint disagreement}.
        \item If checkpoint = terminal block → persistent disagreement (breaks consensus).
        \item If checkpoint = \textbf{parent of terminal block} → miners stay coordinated after latency (Lemma 4, Proposition 5).
        \item \textbf{Implication}: Delayed settlement (e.g., Bitcoin’s 6-block confirmation) aligns with checkpoint lags.
        \item Checkpoint strategies must be designed to absorb latency shocks.
    \end{itemize}
\end{frame}

% Slide 6: Trade-off: Consensus vs. Informational Efficiency
\begin{frame}{Trade-off: Consensus vs. Informational Efficiency}
    \begin{itemize}
        \item \textbf{Full-information checkpoint rules}:
        \begin{itemize}
            \item Use all recent data (e.g., choose chain with highest total balances).
            \item Maximize welfare but fail to achieve consensus (Proposition 6).
            \item Small updates cause large chain shifts → forking incentives.
        \end{itemize}
        \item \textbf{Partial-information checkpoint rules}:
        \begin{itemize}
            \item Ignore some recent info → can sustain consensus.
            \item This may not maximize welfare (Proposition 7).
        \end{itemize}
        \item \textbf{\alert{Fundamental trade-off}}: Cannot have both full informational efficiency and decentralized consensus.
        \item \textbf{Analogous to Grossman-Stiglitz paradox}: fully efficient consensus impossible with strategic agents.
    \end{itemize}
\end{frame}

% Slide 7: Conclusion and Contributions
\begin{frame}{Game-theoretic model of blockchain consensus with rational miners }
  %  \textbf{Key contributions}:
    \begin{enumerate}
        %\item Game-theoretic model of blockchain consensus with rational miners.
        \item Recognizes three attack types under longest-chain rule.
        \item \textbf{Checkpoint strategies} as history-dependent solution:
        \begin{itemize}
            \item Prevent double-spending.
            \item Enhance consensus robustness.
        \end{itemize}
        \item Parent-block checkpoints can be used to manage latency.
        \item There is a  \textbf{Grossman-Stiglitz like trade-off}  between informational efficiency and consensus.
            \end{enumerate}
    
    \vspace{0.3cm}
    \textbf{Implications}: Presence of vulnerabilities is endogenous and  depends on equilibrium strategies. Modifications of protocols (some already in use) can significantly improve security.
\end{frame}

% Slide 8: References
\begin{frame}{References}
    \small
    \begin{itemize}
        \item Biais, B., et al. (2019). \emph{The Blockchain Folk Theorem}. Review of Financial Studies.
        \item Budish, E. (2025). \emph{Trust at Scale: Economic Limits of Cryptocurrencies}. QJE (forthcoming).
        \item Buterin, V., et al. (2020). Ethereum 2.0 finality gadgets.
        \item Leshno, J., et al. (2024). Community-response protocols.
        \item Stewart, A., \& Kokoris-Kogia, E. (2020). GRANDPA: a Byzantine finality gadget.
        \item Garratt, R., \& van Oordt, M. (2023). Double-spending costs.
        \item Chiu, J., \& Koeppl, T. (2022). Miner competition and settlement lags.
        \item {Zetlin-Jones, A., Ebrahimi, Z., Guennewig, M., \& Routledge, B. (2025).} 
\emph{Achieving Consensus on Blockchains}. 
Working paper, November 7, 2025.
    \end{itemize}
\end{frame}

\end{document}

% End slide
\begin{frame}
    \centering
    \Large \textbf{Thank You}
    
    \vspace{0.5cm}
    \small Summary of:\\
    \textbf{Achieving Consensus on Blockchains}\\
    Zetlin-Jones et al., November 2025
\end{frame}

\end{document}